% paper/sections/setup.tex -- P2-01, plan/03-phase2-theory.md.
%
% Standalone section fragment, \input{} by paper/main.tex once P6-01 builds
% the skeleton. Requires (from the main preamble, not redeclared here):
%   amsmath, amssymb, amsthm, mathtools
%   \newtheorem{assumption}{Assumption}
%   \newtheorem{definition}{Definition}
%   \newtheorem{remark}{Remark}
% and a \needshuman{...} macro (defined once main.tex exists; renders its
% argument in a visually distinct box -- see plan/03-phase2-theory.md's
% instruction to mark every non-mechanical proof step this way rather than
% write "it can be shown that").
%
% Every symbol here is mapped to the code identifier that computes or
% estimates it in docs/notation.md -- read that file alongside this one.
%
% Citation keys used below (P6-01 must add these to the project .bib when
% the bibliography is set up; not yet resolvable, no .bib file exists):
%   hoffmann2022chinchilla -- Hoffmann et al. 2022, arXiv:2203.15556
%     (see docs/related_work.md, already vetted there).

\section{Formal setup}
\label{sec:setup}

\subsection{Arms, scales, and the target}

There are $K$ arms (candidate data recipes), indexed $k \in \{1, \dots, K\}$. Each arm
$k$ has a true decision-relevant quantity $\mu_k(s)$ at every scale $s$, where a
\emph{scale} $s = (N, D) \in \mathcal{S}$ is a (parameter count, token count) pair. The
\emph{target scale} $s^\star \in \mathcal{S}$ is the scale at which the decision must
actually be made -- $\mu_k(s^\star)$ is the quantity of interest, and $\mathcal{S}$ may
contain scales both below and (in principle) above $s^\star$, though every scale this
paper actually fits at is below it. The decision-maker never observes $\mu_k(s^\star)$
directly except by paying its full cost (Section~\ref{sec:setup-cost}); the entire
problem is choosing how to spend a compute budget observing $\mu_k(s)$ at scales
$s \ne s^\star$ to infer $\mu_k(s^\star)$ well enough to identify
$k^\star = \arg\max_k \mu_k(s^\star)$.

\subsection{The scaling family}
\label{sec:setup-family}

\begin{assumption}[Parametric scaling family]
\label{asn:family}
There is a known function $g: \Theta \times \mathcal{S} \to \mathbb{R}$, with
$\Theta \subset \mathbb{R}^p$ compact, such that each arm's \emph{in-family} trend is
$g(\theta_k, s)$ for some unknown $\theta_k \in \Theta$. $g$ is continuously
differentiable in $\theta$ on $\Theta$ for every $s \in \mathcal{S}$, and its Jacobian
$J(\theta, s) = \nabla_\theta g(\theta, s)$ is bounded on $\Theta \times \mathcal{S}$:
$\sup_{\theta \in \Theta, s \in \mathcal{S}} \lVert J(\theta, s) \rVert < \infty$.
\end{assumption}

This project implements six concrete choices of $g$ (\texttt{src/pdt/scaling/fitters.py}),
differing in $p$ and functional form:
\texttt{ConstantExtrapolator} ($p = 1$, $g$ does not depend on $s$ -- the single-scale
baseline, Corollary~\ref{cor:single-scale} below), \texttt{LogLinear} ($p = 2$),
\texttt{PowerLawN} and \texttt{PowerLawC} ($p = 3$ each, differing in whether the power
law is indexed by parameter count or compute), \texttt{ChinchillaND} ($p = 5$, the joint
$(N, D)$ form of \citet{hoffmann2022chinchilla}), and \texttt{TwoStepLadder} ($p = 7$).
Assumption~\ref{asn:family} is stated for a generic $g$ because every theorem below
applies to any of the six -- the fitter is a free parameter of the theory, not baked
into a theorem statement, matching how \texttt{src/pdt/theory/bound.py} takes a fitted
\texttt{Extrapolator} instance as an argument rather than hardcoding one.

Fitting $\theta_k$ requires a \emph{design} $S_{\mathrm{fit}} \subset \mathcal{S}$, a
finite multiset of scales at which arm $k$ is actually observed
(Section~\ref{sec:setup-obs}). \texttt{Extrapolator.fit()} refuses a design with
$|S_{\mathrm{fit}}| < p + 1$ (raises \texttt{FitFailure}) -- this is not a numerical
convenience but a real content of the theory, since Theorem~\ref{thm:identifiability}
below is exactly about what $|S_{\mathrm{fit}}|$ and its spacing must satisfy for
$\theta_k$ to be identifiable at all.

\subsection{The observation model}
\label{sec:setup-obs}

\begin{assumption}[Noisy observations]
\label{asn:noise}
At each scale $s \in S_{\mathrm{fit}}$ where arm $k$ is observed, one (or several
independent, indexed $i$) noisy observations are drawn:
\[
  y_{k, i}(s) = \mu_k(s) + \varepsilon_{k, i}(s), \qquad
  \varepsilon_{k, i}(s) \ \text{sub-Gaussian with proxy variance } \sigma^2(s),
\]
independent across $(k, i, s)$, with $\sigma^2(s)$ known or (as in this project's actual
experiments) estimated from repeated observations before the design is chosen.
\end{assumption}

\begin{remark}[The scale-noise assumption is empirically false, and this is a
finding, not a nuisance]
\label{rem:noise-not-monotone}
It is natural to expect $\sigma^2(s)$ to \emph{decrease} in $s$ -- a well-worn
intuition from standard best-arm identification, where more samples (here, larger
models) are usually less noisy per observation, and an earlier draft of this project
assumed exactly that. \textbf{P1-05 measured this directly on DataDecide and found no
such trend}: the seed-variance component of $\sigma^2(s)$ has a median ranging from
about $2.9 \times 10^{-5}$ to $5.9 \times 10^{-5}$ across the 4M-to-1B ladder with no
clear monotonic direction (\texttt{results/p1\_05\_noise.json};
\texttt{docs/decisions.md}, 2026-09-03). Per Rule 2 of \texttt{plan/00-agent-protocol.md}
(``report results that contradict the hypothesis''), Assumption~\ref{asn:noise} is
therefore stated \emph{without} assuming monotonicity in $s$: $\sigma^2(\cdot)$ is a
general, scale-indexed (heteroscedastic) function, not a decreasing one. This matters
for Theorem~\ref{thm:lower-bound}'s allocation program (Part~A): the classical BAI
intuition that bigger scales are doubly attractive (informative \emph{and} eventually
cheap-per-bit) does not hold here on the noise side -- only the bias side
(Section~\ref{sec:setup-misspec}) pulls the allocation toward larger scales, which is a
strictly weaker force. \texttt{P3-02}'s numerical solver for $T^\star$ must therefore
use the \emph{measured} $\sigma^2(s)$ (\texttt{src/pdt/analysis/noise.py}), never an
assumed decreasing curve.
\end{remark}

In practice, $\sigma^2(s)$ has (at least) three independent, separately estimated
components -- seed variance, checkpoint jitter, and eval-sampling noise
(\texttt{src/pdt/analysis/noise.py}: \texttt{seed\_variance}, \texttt{checkpoint\_jitter},
\texttt{eval\_sampling\_noise} / \texttt{eval\_sampling\_noise\_of\_mean}). Which
components apply, and how they combine, depends on what a single ``observation'' $y$
means for a given estimator (a single seed's single-checkpoint score, versus a
seed-averaged score); \texttt{docs/notation.md} states the exact correspondence used
by each results file. Theorem~\ref{thm:main-bound} treats $\sigma^2(s)$ as given and is
agnostic to this decomposition.

\subsection{The cost model}
\label{sec:setup-cost}

\begin{assumption}[Compute cost]
\label{asn:cost}
Observing arm $k$ once at scale $s = (N, D)$ costs $c(s) = 6 N D$ FLOPs (the standard
transformer training-compute approximation; \texttt{Scale.compute} in
\texttt{src/pdt/scaling/base.py}). A design that observes scale $s_i$ a total of $n_i$
times costs $\sum_i n_i \, c(s_i)$, and the total compute spent by a policy up to any
point is $C = \sum_i n_i \, c(s_i)$, summed over everything pulled so far.
\end{assumption}

\subsection{The misspecification model}
\label{sec:setup-misspec}

This is the paper's central modelling choice, and the one requiring the most care.

\begin{assumption}[Bounded misspecification]
\label{asn:misspec}
Each arm's true trend is $\mu_k(s) = g(\theta_k, s) + h_k(s)$, where $h_k$ lies in a
known bounded perturbation class $\mathcal{H}$ -- for concreteness, either
$\sup_{s \in \mathcal{S}} |h_k(s)| \le \eta$ for a known or estimated $\eta \ge 0$, or a
Hölder ball of known smoothness and radius. $g$ is never exactly correct; $\mathcal{H}$
is the family's honest admission of that, not an edge case.
\end{assumption}

\begin{definition}[Population-level in-family projection]
\label{def:projection}
For a design $S_{\mathrm{fit}}$ (formally, a probability measure $\pi$ on
$\mathcal{S}$ supported on the scales observed, weighted by how much of the budget each
receives), the \emph{population-level} (noise-free, infinite-replicate) least-squares
projection of arm $k$'s true trend onto the family is
\[
  \theta_k^{\dagger}(S_{\mathrm{fit}})
  \;=\; \arg\min_{\theta \in \Theta} \ \mathbb{E}_{s \sim \pi}\!\left[
    \big( g(\theta, s) - \mu_k(s) \big)^2 \right].
\]
\end{definition}

\begin{definition}[Extrapolation bias, $\sigma^2_{\mathrm{extrap}}$]
\label{def:sigma2-extrap}
\[
  \sigma^2_{\mathrm{extrap}, k}(S_{\mathrm{fit}})
  \;=\; \Big( g\big(\theta_k^{\dagger}(S_{\mathrm{fit}}), s^\star\big) - \mu_k(s^\star)
  \Big)^2 .
\]
\end{definition}

\begin{remark}[Design-dependence is the mathematical heart of the paper]
\label{rem:design-dependence}
$\sigma^2_{\mathrm{extrap}, k}$ is a property of the pair $(k, S_{\mathrm{fit}})$, not
of $k$ alone -- two different designs project the same true $\mu_k$ onto two different
in-family curves, with two different biases at $s^\star$. Consequently:
\begin{itemize}
  \item \textbf{More replicates at the same scales do not shrink
    $\sigma^2_{\mathrm{extrap}, k}$.} They shrink the \emph{estimation} variance
    $v_k(C)$ (Theorem~\ref{thm:main-bound}) toward $0$, but
    $\theta_k^{\dagger}(S_{\mathrm{fit}})$ -- and hence
    $\sigma^2_{\mathrm{extrap}, k}(S_{\mathrm{fit}})$ -- does not depend on $C$ at all
    for a \emph{fixed design shape} (fixed relative weighting $\pi$ across scales, just
    scaled up). This is exactly the empirical pattern P1-06 measured: across
    $6$ fitters $\times$ $3$ designs $\times$ $11$ tasks, \texttt{v\_hat} shrinks with
    fitting compute while \texttt{sigma2\_extrap\_hat} does not shrink to $0$ at any
    design (\texttt{results/p1\_06\_decomposition.json}; a $\sim$14$\times$ spread
    across tasks at the same fitter and design, \texttt{docs/findings/p1\_06.md}).
  \item \textbf{Changing which scales are in the design changes
    $\sigma^2_{\mathrm{extrap}, k}$}, because it changes $\pi$ and hence
    $\theta_k^{\dagger}$. This is the lever Theorem~\ref{thm:identifiability}'s
    design-dependence corollary formalises, and what P1-06's own (surprising,
    contrary-to-plan) finding is really about: P1-06 found the ratio
    $\sigma^2_{\mathrm{extrap}} / v$ \emph{falls}, not rises, as the design's largest
    scale moves closer to $s^\star$ -- consistent with this remark (a design whose
    scales sit closer to $s^\star$ changes $\pi$ toward the region where $g$'s
    in-family fit is closer to the truth), not with the earlier, wrong prediction that
    $\sigma^2_{\mathrm{extrap}}$ should stay flat in $C$ while only $v$ moves.
\end{itemize}
\end{remark}

\begin{remark}[Estimating $\sigma^2_{\mathrm{extrap}, k}$]
\label{rem:sigma2-extrap-estimator}
Definitions~\ref{def:projection}--\ref{def:sigma2-extrap} are population quantities;
$\mu_k$ and $\theta_k^{\dagger}$ are never observed exactly. P1-06's estimator
(\texttt{bias\_variance\_decomposition} in \texttt{src/pdt/analysis/bootstrap.py})
must separate \emph{three} sources of noise in the naive $\widehat{\mathrm{bias}}_k^2$,
where $\widehat{\mathrm{bias}}_k = \overline{\hat\mu}_k^{(1:B)} - \mu_{k,\mathrm{true}}$
is the mean of $B$ bootstrap predictions of $\mu_k(s^\star)$ minus the measured ground
truth:
\begin{enumerate}
  \item \emph{Monte Carlo noise of the bootstrap mean.} Given the observed dataset the
    $B$ replicates are independent, so $\overline{\hat\mu}^{(1:B)}$ has variance
    $\hat v_k / B$ around the bootstrap population mean. This vanishes as
    $B \to \infty$.
  \item \emph{Sampling noise of the original fit.} As $B \to \infty$ the bootstrap mean
    converges to a functional of the \emph{observed} dataset -- essentially the original
    fit's own prediction $\hat\mu_k^{\mathrm{orig}}(s^\star)$ -- not to
    $\mathbb{E}[\hat\mu_k^{\mathrm{orig}}(s^\star)]$. So $\widehat{\mathrm{bias}}_k^2$
    keeps the original estimator's sampling variance
    $\mathrm{Var}(\hat\mu_k^{\mathrm{orig}}(s^\star))$ \emph{for every $B$}. (Example: the
    mean of $n$ seeds with true target exactly equal to the population mean has squared
    structural bias $0$, yet $\mathbb{E}[\widehat{\mathrm{bias}}^2] \to \sigma^2/n$ as
    $B \to \infty$.) Dividing $\hat v_k$ by $B$ therefore does \emph{not} remove it.
  \item \emph{Noise in the ground truth} $\sigma^2_{\mathrm{target}, k}$
    (Assumption~\ref{asn:noise} at $s^\star$).
\end{enumerate}
The bootstrap replicate variance $\hat v_k$ estimates (2) but, for $n$-out-of-$n$
resampling of $n$ seeds, is the \emph{plug-in} variance: a factor $(n-1)/n$ below the
unbiased $\mathrm{Var}(\hat\mu_k^{\mathrm{orig}})$ (exactly, for a sample mean). The
implemented, approximately-unbiased estimator of the squared \emph{structural} bias is
therefore
\[
  \widehat{\mathrm{bias}^2_k}\;=\;\widehat{\mathrm{bias}}_k^2
    - \frac{n}{n-1}\,\hat v_k \;-\; \frac{\hat v_k}{B} \;-\; \sigma^2_{\mathrm{target}, k},
  \qquad
  \widehat{\sigma^2_{\mathrm{extrap}, k}}\;=\;\max\!\big(0,\ \widehat{\mathrm{bias}^2_k}\big),
\]
with $n/(n-1) = 3/2$ for the three DataDecide seeds (and $1$ for a parametric bootstrap,
which draws from an assumed noise model rather than resampling $n$ observations).
Status of each claim: the \emph{unclipped} $\widehat{\mathrm{bias}^2_k}$ is exactly
unbiased for a sample mean and a first-order (delta-method) approximation for the smooth
fitters used here; it is an approximation at $n = 3$ and for non-smooth or
boundary-pinned fits, not a proof. The \emph{clipped} $\widehat{\sigma^2_{\mathrm{extrap}, k}}$
is a nonnegative \textbf{heuristic}: $\mathbb{E}[\max(0, X)] > \mathbb{E}[X]$, so it is
biased \emph{upward} when the true bias is small, and any average of clipped values
overstates the mean squared bias (average the unclipped values to estimate it). The
population quantity is nonnegative; the estimate is negative about half the time when the
true bias is near zero, which is why the clip exists at all.
\end{remark}

\subsection{The policy class and correctness criterion}
\label{sec:setup-policy}

A policy is a (possibly randomised, adaptive) rule that, at each step, either selects
an arm-scale pair $(k, s) \in \{1, \dots, K\} \times \mathcal{S}$ to observe next or
stops. Formally: a stopping time $\tau$ with respect to the filtration generated by the
observations so far, together with a recommendation $\hat{k} \in \{1, \dots, K\}$ that
is $\mathcal{F}_\tau$-measurable. Let
$k^\star = \arg\max_k \mu_k(s^\star)$ (assumed unique unless stated otherwise) and
$\Delta_k = \mu_{k^\star}(s^\star) - \mu_k(s^\star) \ge 0$ for $k \ne k^\star$.

\begin{definition}[$\delta$-correctness]
\label{def:delta-correct}
A policy is $\delta$-correct on an instance class $\mathcal{I}$ if
$\mathbb{P}[\hat{k} \ne k^\star] \le \delta$ for every instance in $\mathcal{I}$.
\end{definition}

Section~\ref{sec:setup-misspec}'s misspecification model is folded into
$\mathcal{I}$ via $\mathcal{H}$: an instance is a full specification
$(\theta_1, \dots, \theta_K, h_1, \dots, h_K)$ with each $h_k \in \mathcal{H}$, and
$\delta$-correctness (Definition~\ref{def:delta-correct}) is required uniformly over
every such instance, not just over the noise given fixed, correctly-specified
$\theta_k$'s -- this is what makes Theorem~\ref{thm:main-bound}'s bound depend on
$\sigma^2_{\mathrm{extrap}, k}$ and not only on estimation variance, and what makes
Theorem~\ref{thm:lower-bound}'s impossibility result (Part~B) possible at all.

\bigskip
\noindent\emph{Every symbol introduced above is mapped to its computing code identifier
in \texttt{docs/notation.md}.}
